\documentclass[../../main.tex]{subfiles}% 注意这里的文件路径不能用 ./main.tex ,否则用latexmk编译子文件会报错
\graphicspath{{\subfix{./image/}}{\subfix{../../image/}}} % 指定图片引用路径,后续可以直接使用图片文件名
% 如果图片存放在上级目录,则使用 ../../image/ 作为备用路径

% 例如:
% \begin{figure}[H]
% \centering
% \includegraphics[width=0.3\textwidth]{图.png}
% \caption{}
% \label{figure:图}
% \end{figure}
% 注意:上述\label{}一定要放在\caption{}之后,否则引用图片序号会只会显示??.

\begin{document}

\section{Hilbert-Schmidt定理}

\begin{definition}
设$H$是Hilbert空间,$A\in\mathscr{L}(H)$，其\textbf{共轭算子}$A^*$由下式定义：
\begin{align}
(Ax,y)=(x,A^*y)\quad(\forall x,y\in H).\label{eq:conj-op-def}
\end{align}
称$A$是\textbf{对称的}，若
\begin{align}
(Ax,y)=(x,Ay)\quad(\forall x,y\in H).\label{eq:sym-op-def}
\end{align}
比较\eqref{eq:conj-op-def}式和\eqref{eq:sym-op-def}式可见，$A$是对称算子当且仅当$A=A^*$。因此对称算子也称为\textbf{自共轭算子}或\textbf{自伴算子}.
\end{definition}

\begin{proposition}[proposition:共轭算子的基本性质]\label{proposition:共轭算子的基本性质}
设$H$是Hilbert空间,$A,B\in\mathscr{L}(H),\alpha\in\mathbb{C}$，则有
\begin{align*}
(A+B)^*&=A^*+B^*,\quad (\alpha A)^*=\overline{\alpha}A^*,\\
A^{**}&=A,\quad (AB)^*=B^*A^*,\\
\sigma(A^*)&=\overline{\sigma(A)}=\{\overline{\lambda}\mid\lambda\in\sigma(A)\}.
\end{align*}
\end{proposition}
\begin{proof}
由共轭算子定义直接验证可得。
\end{proof}

\begin{examples}
在$\mathbb{R}^n$上，若$A$是对称矩阵，则$A$是对称算子；在$\mathbb{C}^n$上，若$A$是Hermite矩阵，则$A$是对称算子。
一般地：
\begin{enumerate}[(1)]
\item $A$是$\mathbb{R}^n$上的矩阵，$A^*=A^\mathrm{T}$，其中$A^\mathrm{T}$为$A$的转置矩阵；

\item $A$是$\mathbb{C}^n$上的矩阵，$A^*=\overline{A^\mathrm{T}}$，其中$\overline{A^\mathrm{T}}$为$A$的共轭转置矩阵。
\end{enumerate}
\end{examples}

\begin{examples}
在实的$L^2(\Omega,\mathscr{B},\mu)$上，设$K\in L^\infty(\Omega\times\Omega,\mathrm{d}\mu)$，并且$K(x,y)=K(y,x)$，则算子
\begin{align*}
A:u(x)\mapsto\int_{\Omega}K(x,y)u(y)\mathrm{d}\mu(y)
\end{align*}
是$L^2(\Omega,\mathscr{B},\mu)$上的对称算子。
\end{examples}

\begin{examples}
设$H$是Hilbert空间，$M$是其闭线性子空间，则由$H$到$M$上的投影算子$P_M$是对称算子。
\end{examples}
\begin{proof}
由正交分解定理，$\forall x,y\in H$，有分解
\begin{align*}
x=x_M+x_{M^\perp}\quad(x_M\in M,x_{M^\perp}\in M^\perp),\\
y=y_M+y_{M^\perp}\quad(y_M\in M,y_{M^\perp}\in M^\perp).
\end{align*}
由$P_M$的定义，$x_M=P_Mx,\,y_M=P_My$，因此
\begin{align*}
(P_Mx,y)=(x_M,y_M+y_{M^\perp})=(x_M,y_M)=(x,P_My).
\end{align*}
\end{proof}

\begin{proposition}
设$H$是Hilbert空间,关于$H$上的对称算子，有下列基本性质：
\begin{enumerate}[(1)]
\item $A$对称$\iff(Ax,x)\in\mathbb{R}\ (\forall x\in H)$.
\item 若$A$对称，则$\sigma(A)\subseteq\mathbb{R}$，且
\begin{align*}
\|(\lambda I-A)^{-1}x\|\leqslant\frac{1}{|\mathrm{Im}\lambda|}\|x\|\quad(\forall x\in H,\forall\lambda\in\mathbb{C},\mathrm{Im}\lambda\neq0).
\end{align*}
\item 设$H_1$是$H$的闭不变子空间，$A$是$H$上的对称算子，则$A|_{H_1}$也是$H_1$上的对称算子。
\item 若$A$对称，$\lambda,\lambda'\in\sigma_p(A),\lambda\neq\lambda'$，则$N(\lambda I-A)\perp N(\lambda'I-A)$.
\item 若$A$对称，则$\displaystyle\sup_{\|x\|=1}|(Ax,x)|=\|A\|$.
\end{enumerate}
\end{proposition}
\begin{proof}
\begin{enumerate}[(1)]
\item 令$a(x,y)\triangleq(Ax,y)\ (\forall x,y\in H)$，则$a(\cdot,\cdot)$是$H$上的共轭双线性函数，$(Ax,x)$是由$a(\cdot,\cdot)$诱导的二次型。
由定义，
\begin{align}
A\text{对称}\iff a(x,y)=\overline{a(y,x)}\quad(\forall x,y\in H).\label{eq:sym-quad}
\end{align}
又共轭双线性函数性质：
\begin{align}
a(x,y)=\overline{a(y,x)}\quad(\forall x,y\in H)\iff(Ax,x)\in\mathbb{R}\quad(\forall x\in H).\label{eq:real-quad}
\end{align}
联合\eqref{eq:sym-quad}与\eqref{eq:real-quad}即得结论。

\item 设$\lambda=\mu+\mathrm{i}\nu,\nu\neq0,\mu,\nu\in\mathbb{R}$，由$A$的对称性，
\begin{align*}
\|(\lambda I-A)x\|^2=\|(\mu I-A)x\|^2+|\nu|^2\cdot\|x\|^2
\geqslant|\nu|^2\cdot\|x\|^2\quad(\forall x\in H).
\end{align*}
此外，
\begin{align*}
R(\lambda I-A)^\perp=N(\overline{\lambda}I-A^*)=N(\overline{\lambda}I-A),
\end{align*}
结合上式可得$N(\overline{\lambda}I-A)=\{\theta\}\ (\mathrm{Im}\lambda\neq0)$，故$R(\lambda I-A)=H$。

\item 由对称算子定义直接验证。

\item 若$x\in N(\lambda I-A),\,x'\in N(\lambda'I-A)$，则
\begin{align*}
\lambda(x,x')=(Ax,x')=(x,Ax')=\lambda'(x,x').
\end{align*}
由$\lambda\neq\lambda'$，推出$(x,x')=0$。

\item 记$\displaystyle c\triangleq\sup_{\|x\|=1}|(Ax,x)|$，显然$c\leqslant\|A\|$。下证$c\geqslant\|A\|$。
由$A$的对称性和\hyperref[proposition:平行四边形等式--泛函分析]{平行四边形法则}，
\begin{align*}
\mathrm{Re}(Ax,y)&=\frac{1}{4}\left[(A(x+y),x+y)-(A(x-y),x-y)\right]\\
&\leqslant\frac{c}{4}\left(\|x+y\|^2+\|x-y\|^2\right)\leqslant c,
\end{align*}
其中$x,y\in H$满足$\|x\|=\|y\|=1$。
取$\alpha\in\mathbb{C}(|\alpha|=1)$，使得$\alpha(Ax,y)=|(Ax,y)|$，则
\begin{align*}
|(Ax,y)|=(Ax,\overline{\alpha}y)=\mathrm{Re}(Ax,\overline{\alpha}y)\leqslant c.
\end{align*}
故$\|A\|\leqslant c$。
\end{enumerate}
\end{proof}

\begin{theorem}\label{theorem:对称算子的算子范数等于某个元的内积}
若$A$是对称紧算子，则存在$x_0\in H,\|x_0\|=1$，使得
\begin{align*}
|(Ax_0,x_0)|=\sup_{\|x\|=1}|(Ax,x)|,
\end{align*}
且满足
\begin{align}
Ax_0=\lambda x_0,\label{eq:sym-compact-eig}
\end{align}
其中$|\lambda|=|(Ax_0,x_0)|$.
\end{theorem}
\begin{proof}
记$S_1$为$H$的单位球面，不妨设
\begin{align}
\sup_{x\in S_1}|(Ax,x)|=\sup_{x\in S_1}(Ax,x).\label{eq:sym-compact-sup}
\end{align}
（否则用$-A$代替$A$）。取$\displaystyle\lambda\triangleq\sup_{x\in S_1}(Ax,x)$，定义函数$f(x)=(Ax,x)\ (\forall x\in S_1)$。
取序列$\{x_n\}\subset S_1$满足$f(x_n)\to\lambda$，因$\|x_n\|=1$，故存在弱收敛子列，不妨仍记为$\{x_n\}$，设$x_n\rightharpoonup x_0$，由闭球$\overline{B(\theta,1)}$的弱闭性得$\|x_0\|\leqslant1$。
由$A$的紧性得$Ax_n\to Ax_0\ (n\to\infty)$，故$f(x_n)\to(Ax_0,x_0)$，即$(Ax_0,x_0)=\lambda$。

下证$\|x_0\|=1$，反证：若$\|x_0\|<1$，则
\begin{align*}
(Ax_0,x_0)\leqslant\|A\|\cdot\|x_0\|^2=\sup_{x\in S_1}(Ax,x)\|x_0\|^2<\lambda,
\end{align*}
矛盾，故$\exists x_0\in S_1$，使得
\begin{align*}
(Ax_0,x_0)=\sup_{x\in S_1}(Ax,x).
\end{align*}

最后证\eqref{eq:sym-compact-eig}。$\forall y\in H$，对充分小的$t$，定义
\begin{align*}
\varphi_y(t)\triangleq\frac{(A(x_0+ty),x_0+ty)}{(x_0+ty,x_0+ty)}.
\end{align*}
$t=0$使$\varphi_y(t)$取极大值，故$\varphi_y'(0)=0$，计算得
\begin{align*}
\mathrm{Re}(y,Ax_0-\lambda x_0)=0,
\end{align*}
由$y$的任意性，得$Ax_0=\lambda x_0$。
\end{proof}

\begin{theorem}[Hilbert-Schmidt定理]\label{theorem:Hilbert-Schmidt定理}
若$A$是Hilbert空间$H$上的对称紧算子，则至多有可数个非零、仅能以$0$为聚点的实数$\{\lambda_i\}$，它们是$A$的特征值，并对应一组正交规范基$\{e_i\}$，使得对$\forall x\in H$,有
\begin{align}
x=\sum\limits_{i}(x,e_i)e_i,\quad 
Ax=\sum\limits_{i}\lambda_i(x,e_i)e_i.\label{eq:hs-exp-ax}
\end{align}
\end{theorem}
\begin{remark}
\begin{enumerate}[(1)]
\item 将特征值按绝对值递减顺序编号，约定重数为几则重复编号，即
\begin{align*}
|\lambda_1|\geqslant|\lambda_2|\geqslant\cdots\geqslant|\lambda_n|\geqslant|\lambda_{n+1}|\geqslant\cdots,
\end{align*}
于是
\begin{align*}
A=\sum\limits_{i=1}^{\infty}\lambda_i e_i\otimes e_i,
\end{align*}
且
\begin{align*}
\left\|A-\sum\limits_{i=1}^{n}\lambda_i e_i\otimes e_i\right\|\leqslant|\lambda_{n+1}|\to0\quad(n\to\infty).
\end{align*}
事实上，$\forall x\in H$，
\begin{align*}
\left\|Ax-\sum\limits_{i=1}^{n}\lambda_i(x,e_i)e_i\right\|&=\left\|\sum\limits_{i=n+1}^{\infty}\lambda_i(x,e_i)e_i\right\|
=\left(\sum\limits_{i=n+1}^{\infty}\lambda_i^2|(x,e_i)|^2\right)^{\frac{1}{2}}\\
&\leqslant|\lambda_{n+1}|\left(\sum\limits_{i=n+1}^{\infty}|(x,e_i)|^2\right)^{\frac{1}{2}}\leqslant|\lambda_{n+1}|\cdot\|x\|.
\end{align*}
\item \hyperref[theorem:Hilbert-Schmidt定理]{Hilbert-Schmidt定理}表明:对称紧算子可对角化，其特征值具有极值性质：
\begin{align*}
|\lambda_n|=\sup\left\{|(Ax,x)|\mid x\perp\mathrm{span}\{e_1,e_2,\cdots,e_{n-1}\},\|x\|=1\right\}\quad(n=1,2,\cdots),
\end{align*}
其中$e_1,e_2,\cdots,e_{n-1}$是对应$\lambda_1,\lambda_2,\cdots,\lambda_{n-1}$的特征元。

可按正负值排列特征值：
\begin{align*}
\lambda_1^+\geqslant\lambda_2^+\geqslant\cdots\geqslant0,\\
\lambda_1^-\leqslant\lambda_2^-\leqslant\cdots<0.
\end{align*}
\end{enumerate}
\end{remark}
\begin{proof}
$\forall\lambda\in\sigma_p(A)\setminus\{0\}$，设$N(\lambda I-A)$的正交规范基为$\{e_i^{(\lambda)}\}_{i=1}^{m(\lambda)}$，其中$m(\lambda)\triangleq\dim N(\lambda I-A)<\infty$，称为$\lambda$的重数。
若$0\in\sigma_p(A)$，设$N(A)$的正交规范基为$\{e_i^{(0)}\}$。

令
\begin{align*}
\{e_i'\}\triangleq\bigcup\limits_{\lambda\in\sigma_p(A)\setminus\{0\}}\{e_i^{(\lambda)}\}_{i=1}^{m(\lambda)},
\end{align*}
\begin{align*}
\{e_i\}\triangleq
\begin{cases}
\{e_i'\},&0\notin\sigma_p(A),\\
\{e_i'\}\cup\{e_i^{(0)}\},&0\in\sigma_p(A).
\end{cases}
\end{align*}
记$M\triangleq\mathrm{span}\{e_i\}$，则在$M$上$A$满足\eqref{eq:hs-exp-ax}式。

下证$\overline{M}=H$，反证：若$M^\perp\neq\{\theta\}$，记$\widetilde{A}\triangleq A|_{M^\perp}$，则$\widetilde{A}$无特征值，故$\widetilde{A}\neq0$。
由\refthe{theorem:对称算子的算子范数等于某个元的内积}知
\begin{align*}
\|\widetilde{A}\|=\sup_{\substack{x\in M^\perp\\\|x\|=1}}|(\widetilde{A}x,x)|=0,
\end{align*}
即$\widetilde{A}=0$，矛盾。故$\{e_i\}$构成$H$的正交规范基。
\end{proof}

\begin{theorem}[极小极大刻画]
设$A$是Hilbert空间$H$上的对称紧算子，对应特征值为
\begin{align*}
\lambda_1^+\geqslant\lambda_2^+\geqslant\cdots\geqslant0,\\
\lambda_1^-\leqslant\lambda_2^-\leqslant\cdots<0.
\end{align*}
则
\begin{align}
\lambda_n^+=\inf_{E_{n-1}}\sup_{\substack{x\in E_{n-1}^\perp\\x\neq\theta}}\frac{(Ax,x)}{(x,x)},\label{eq:minimax-pos}\\
\lambda_n^-=\sup_{E_{n-1}}\inf_{\substack{x\in E_{n-1}^\perp\\x\neq\theta}}\frac{(Ax,x)}{(x,x)},\label{eq:minimax-neg}
\end{align}
其中$E_{n-1}$是$H$的任意$n-1$维闭线性子空间。
\end{theorem}
\begin{proof}
只需证\eqref{eq:minimax-pos}，用$-A$代替$A$可得\eqref{eq:minimax-neg}。
若$\displaystyle x=\sum\limits_{j}a_j^+e_j^++\sum\limits_{j}a_j^-e_j^-$，则
\begin{align*}
\frac{(Ax,x)}{(x,x)}=\frac{\sum\limits_{j}\lambda_j^+|a_j^+|^2+\sum\limits_{j}\lambda_j^-|a_j^-|^2}{\sum\limits_{j}|a_j^+|^2+\sum\limits_{j}|a_j^-|^2}.
\end{align*}
记\eqref{eq:minimax-pos}右端为$\mu_n$，分两步证$\lambda_n^+=\mu_n$：
\begin{enumerate}[(1)]
\item $\lambda_n^+\leqslant\mu_n$：
$\forall E_{n-1}$，在$\mathrm{span}\{e_1^+,e_2^+,\cdots,e_n^+\}$中存在非零向量$x_n\perp E_{n-1}$，于是
\begin{align*}
\sup_{\substack{x\perp E_{n-1}\\x\neq\theta}}\frac{(Ax,x)}{(x,x)}\geqslant\frac{(Ax_n,x_n)}{(x_n,x_n)}=\frac{\sum\limits_{j=1}^{n}\lambda_j^+|a_j^+|^2}{\sum\limits_{j=1}^{n}|a_j^+|^2}\geqslant\lambda_n^+,
\end{align*}
故$\lambda_n^+\leqslant\mu_n$。
\item $\lambda_n^+\geqslant\mu_n$：
取$E_{n-1}=\mathrm{span}\{e_1^+,e_2^+,\cdots,e_{n-1}^+\}$，则
\begin{align*}
\lambda_n^+=\sup_{\substack{x\perp E_{n-1}\\x\neq\theta}}\frac{(Ax,x)}{(x,x)},
\end{align*}
故$\lambda_n^+\geqslant\mu_n$。
\end{enumerate}
\end{proof}

\begin{corollary}
若Hilbert空间$H$上的两个对称紧算子$A,B$满足$A\leqslant B$，即
\begin{align*}
(Ax,x)\leqslant(Bx,x)\quad(\forall x\in H),
\end{align*}
对应特征值为
\begin{gather*}
\lambda _{1}^{+}\left( A \right) \geqslant \lambda _{2}^{+}\left( A \right) \geqslant \cdots \geqslant 0,\quad \lambda _{1}^{-}\left( A \right) \leqslant \lambda _{2}^{-}\left( A \right) \leqslant \cdots <0;
\\
\lambda _{1}^{+}\left( B \right) \geqslant \lambda _{2}^{+}\left( B \right) \geqslant \cdots \geqslant 0,\quad \lambda _{1}^{-}\left( B \right) \leqslant \lambda _{2}^{-}\left( B \right) \leqslant \cdots <0.
\end{gather*}
则
\begin{align*}
\lambda_j^+(A)\leqslant\lambda_j^+(B)\quad(j=1,2,\cdots).
\end{align*}
\end{corollary}


(本节各题中, $H$ 均指复Hilbert空间)

\begin{example}
设 $A \in \mathscr{L}(H)$，求证：$A+A^*$，$AA^*$，$A^*A$ 都是对称算子，并且
\begin{align*}
\|AA^*\| = \|A^*A\| = \|A\|^2.
\end{align*}
\end{example}

\begin{example}
设 $A \in \mathscr{L}(H)$，满足 $(Ax,x) \geqslant 0(\forall x \in H)$，且
\begin{align*}
(Ax,x) = 0 \iff x = \theta,
\end{align*}
求证：
\begin{align*}
\|Ax\|^2 \leqslant \|A\|(Ax,x) \quad (\forall x \in H).
\end{align*}
\end{example}

\begin{example}
设 $A$ 是 $H$ 上的有界对称算子，令
\begin{align*}
m(A) \triangleq \inf\limits_{\|x\|=1} (Ax,x), \quad M(A) \triangleq \sup\limits_{\|x\|=1} (Ax,x).
\end{align*}
求证：
\begin{enumerate}[(1)]
\item $\sigma(A) \subseteq [m(A),M(A)]$，且 $m(A), M(A) \in \sigma(A)$。

进一步假设 $A$ 是 $H$ 上的对称紧算子，求证：
\item 若 $m(A) \neq 0$，则 $m(A) \in \sigma_p(A)$；
\item 若 $M(A) \neq 0$，则 $M(A) \in \sigma_p(A)$。
\end{enumerate}
\end{example}

\begin{example}
设 $A$ 是对称紧算子，求证：
\begin{enumerate}[(1)]
\item 若 $A$ 非零，则 $A$ 至少有一个不等于零的特征值；
\item 若 $M$ 是 $A$ 的非零不变子空间，则 $M$ 上必含有 $A$ 的特征元。
\end{enumerate}
\end{example}

\begin{proposition}\label{proposition:正交投影算子的充要条件}
设$H$是Hilbert空间,则
\begin{enumerate}[(1)]
\item\label{proposition:正交投影算子的充要条件-1} $P \in \mathscr{L}(H)$ 是一个正交投影算子当且仅当$P$是对称且幂等的.

\item\label{proposition:正交投影算子的充要条件-2} $P \in \mathscr{L}(H)$ 是一个正交投影算子当且仅当$(Px,x) = \|Px\|^2(\forall x \in H)$.
\end{enumerate}
特别地,若 $P \in \mathscr{L}(H)$ 是一个正交投影算子,则$P$就是$H$到$\text{Im}P$上的正交投影算子.
\end{proposition}
\begin{proof}
\begin{enumerate}[(1)]
\item {\heiti 必要性:}设$P$是$H$到$M$的正交投影算子.对$\forall x,y\in H$, 则
\begin{gather*}
y=Py+\left( y-Py \right),\text{其中}Py\in M,\quad y-Py\in M^\perp;
\\
x=Px+\left( x-Px \right),\text{其中}Px\in M,\quad x-Px\in M^\perp.
\end{gather*}
又因为$P$在$M$是恒等算子,所以$P^2 x=P(Px)=Px$,故$P$是幂等的.并且
\begin{align*}
\left\langle Px,y \right\rangle &=\left\langle Px,Py \right\rangle +\left\langle Px,y-Py \right\rangle =\left\langle Px,Py \right\rangle ; \\
\left\langle x,Py \right\rangle &=\left\langle Px,Py \right\rangle +\left\langle x-Px,Py \right\rangle =\left\langle Px,Py \right\rangle .
\end{align*}
故$\left\langle Px,y \right\rangle =\left\langle x,Py \right\rangle$, 即$P$是对称的.

{\heiti 充分性:}对$\forall x\in H$, 有
\begin{align*}
x=Px+\left( x-Px \right).
\end{align*}
显然$Px\in \mathrm{Im}P$, 由幂等性知
\begin{align*}
P\left( x-Px \right) =Px-P^2x=0\Longrightarrow x-Px\in \mathrm{Ker}P.
\end{align*}
故$H=\mathrm{Im}P+\mathrm{Ker}P$. 设$Px_0\in \mathrm{Im}P\cap \mathrm{Ker}P$, 则
\begin{align*}
Px_0=P^2x_0=P\left( Px_0 \right) =0.
\end{align*}
故$\mathrm{Im}P\cap \mathrm{Ker}P=\left\{ 0 \right\}$, 即
\begin{align*}
H=\mathrm{Im}P\oplus \mathrm{Ker}P.
\end{align*}
故对$\forall x\in H$, 存在唯一的分解
\begin{align*}
x=Px+\left( x-Px \right),\text{其中}Px\in \mathrm{Im}P,x-Px\in \mathrm{Ker}P.
\end{align*}
再对$\forall Py\in \mathrm{Im}P$, 由幂等性和对称性知
\begin{align*}
\left\langle x-Px,Py \right\rangle =\left\langle P\left( x-Px \right),y \right\rangle =\left\langle Px-P^2x,y \right\rangle =\left\langle 0,y \right\rangle =0.
\end{align*}
因此$x-Px\in \left( \mathrm{Im}P \right) ^{\perp}$. 故由\hyperref[theorem:正交投影算子定义]{正交投影算子定义}知$P$就是$H$到$\mathrm{Im}P$的正交投影算子.

\item {\heiti 必要性:}设$P$是$H$到$M$的正交投影算子.对$\forall x\in H$,有$x=Px+\left( x-Px \right) $,其中$Px\in M,x-Px\in M^{\bot}$.从而
\begin{align*}
\left< Px,x \right> =\left< Px,Px \right> +\left< Px,x-Px \right> =\left\| Px \right\| ^2.
\end{align*}


{\heiti 充分性:}先证对称性。对$\forall z\in H$, 有
\begin{align*}
\left\langle Pz,z-Pz \right\rangle =\left\langle Pz,z \right\rangle -\left\| Pz \right\| ^2=0.
\end{align*}
对$\forall x,y\in H$, 分别将$z=x+y,z=x-y,z=x+\mathrm{i}y,z=x-\mathrm{i}y$代入上式得
\begin{align}
\left\langle P\left( x+y \right) ,\left( x+y \right) -P\left( x+y \right) \right\rangle &=\left\langle Px,x-Px \right\rangle +\left\langle Py,y-Py \right\rangle +\left\langle Px,y-Py \right\rangle +\left\langle Py,x-Px \right\rangle =0; \label{eq:orth_proj_sym_31_1} \\
\left\langle P\left( x-y \right) ,\left( x-y \right) -P\left( x-y \right) \right\rangle &=\left\langle Px,x-Px \right\rangle +\left\langle Py,y-Py \right\rangle -\left\langle Px,y-Py \right\rangle -\left\langle Py,x-Px \right\rangle =0; \label{eq:orth_proj_sym_31_2} \\
\left\langle P\left( x+\mathrm{i}y \right) ,\left( x+\mathrm{i}y \right) -P\left( x+\mathrm{i}y \right) \right\rangle &=\left\langle Px,x-Px \right\rangle +\left\langle Py,y-Py \right\rangle -\mathrm{i}\left\langle Px,y-Py \right\rangle +\mathrm{i}\left\langle Py,x-Px \right\rangle =0; \label{eq:orth_proj_sym_31_3} \\
\left\langle P\left( x-\mathrm{i}y \right) ,\left( x-\mathrm{i}y \right) -P\left( x-\mathrm{i}y \right) \right\rangle &=\left\langle Px,x-Px \right\rangle +\left\langle Py,y-Py \right\rangle +\mathrm{i}\left\langle Px,y-Py \right\rangle -\mathrm{i}\left\langle Py,x-Px \right\rangle =0. \label{eq:orth_proj_sym_31_4}
\end{align}
由\eqref{eq:orth_proj_sym_31_1}\eqref{eq:orth_proj_sym_31_2}式相减, \eqref{eq:orth_proj_sym_31_3}\eqref{eq:orth_proj_sym_31_4}式相减得
\begin{align*}
\left\langle Px,y-Py \right\rangle +\left\langle Py,x-Px \right\rangle =0,\quad -\left\langle Px,y-Py \right\rangle +\left\langle Py,x-Px \right\rangle =0.
\end{align*}
上面两式再相加、相减得
\begin{align*}
\left\langle Px,y-Py \right\rangle =\left\langle Py,x-Px \right\rangle =0\Longrightarrow \left\langle Px,y \right\rangle =\left\langle Px,Py \right\rangle ,\quad \left\langle Py,x \right\rangle =\left\langle Py,Px \right\rangle .
\end{align*}
进而
\begin{align*}
\left\langle Px,y \right\rangle =\left\langle Px,Py \right\rangle =\overline{\left\langle Py,Px \right\rangle }=\overline{\left\langle Py,x \right\rangle }=\left\langle x,Py \right\rangle .
\end{align*}
故$P$是对称的。

再证幂等性。对$\forall x\in H$, 由对称性和充分性假设有
\begin{align*}
\left\langle Px,x \right\rangle =\left\| Px \right\| ^2=\left\langle Px,Px \right\rangle =\left\langle P^2x,x \right\rangle \Longrightarrow \left\langle \left( P-P^2 \right) x,x \right\rangle =0.
\end{align*}
令$S\triangleq P-P^2$, 对$\forall y\in H$, 由上式可知
\begin{align*}
\left\langle Sx,y \right\rangle =\frac{1}{4}\left[ \left\langle S\left( x+y \right) ,x+y \right\rangle -\left\langle S\left( x-y \right) ,x-y \right\rangle +\mathrm{i}\left\langle x+\mathrm{i}y,S\left( x+\mathrm{i}y \right) \right\rangle -\mathrm{i}\left\langle x-\mathrm{i}y,S\left( x-\mathrm{i}y \right) \right\rangle \right] =0.
\end{align*}
取$y=Sx$, 即得
\begin{align*}
\left\| \left( P-P^2 \right) x \right\| ^2=\left\| Sx \right\| ^2=\left\langle Sx,Sx \right\rangle =0\Longrightarrow \left( P-P^2 \right) x=0.
\end{align*}
再由$x$的任意性知$P=P^2$. 故$P$是幂等的。
综上，由(1)可知$P$是一个正交投影算子。
\end{enumerate}
\end{proof}

\begin{proposition}
设 $L, M$ 是Hilbert空间$H$ 上的闭线性子空间,$P_M,P_L$分别是$H$到$M,L$的正交投影算子,则
\begin{enumerate}[(1)]
\item $L \perp M \iff P_L P_M = 0$;
\item $L = M^\perp \iff P_L + P_M = I$ (恒同算子);
\item $P_L P_M = P_{L \cap M} \iff P_L P_M = P_M P_L$.
\end{enumerate}
\end{proposition}
\begin{proof}
\begin{enumerate}[(1)]
\item {\heiti 必要性:} 若 $L\perp M$，则对任意 $x\in H$，有 $P_M x\in M\subseteq L^\perp$，从而 $P_L(P_M x)=0$，即 $P_LP_M=0$.

{\heiti 充分性:} 若 $P_LP_M=0$，则对任意 $x\in M$，$x=P_Mx$，所以 $P_L x=P_LP_Mx=0$，故 $x\in \ker P_L = L^\perp$，即 $M\subseteq L^\perp$,故$L\bot M.$

\item {\heiti 必要性:} 若 $L=M^\perp$，则 $H=L\oplus M$，任意 $x\in H$ 可唯一分解为 $x=x_L+x_M$，其中 $x_L\in L,\ x_M\in M$，且正交分解给出 $P_Lx=x_L,\ P_Mx=x_M$，故 $(P_L+P_M)x=x_L+x_M=x$，所以 $P_L+P_M=I$.

{\heiti 充分性:} 若 $P_L+P_M=I$，则对任意 $x\in M^\bot$，$P_Mx=0$，于是 $x=P_Lx+P_Mx=P_Lx\in L$，所以$M^\bot \subseteq L$. 同理，对任意 $x\in L$，$P_Lx=x$，得 $x=x+P_Mx$，即 $P_Mx=0$，所以 $x\in M^\perp$，即 $L\subseteq M^\perp$. 故 $M^\perp=L$.

\item {\heiti 必要性:}由\rrefpro{proposition:正交投影算子的充要条件}{proposition:正交投影算子的充要条件-1}和\hyperref[proposition:共轭算子的基本性质]{共轭算子的基本性质}知
\begin{align*}
P_LP_M=P_{L\cap M}=P_{L\cap M}^{*}=\left( P_LP_M \right) ^*=P_{M}^{*}P_{L}^{*}=P_MP_L.
\end{align*}

{\heiti 充分性:}由\rrefpro{proposition:正交投影算子的充要条件}{proposition:正交投影算子的充要条件-1}和\hyperref[proposition:共轭算子的基本性质]{共轭算子的基本性质}知
\begin{gather*}
\left( P_LP_M \right) ^*=P_{M}^{*}P_{L}^{*}=P_MP_L=P_LP_M,
\\
\left( P_LP_M \right) ^2=\left( P_LP_M \right) \left( P_LP_M \right) =P_LP_LP_MP_M=P_{L}^{2}P_{M}^{2}=P_LP_M.
\end{gather*}
再结合\rrefpro{proposition:正交投影算子的充要条件}{proposition:正交投影算子的充要条件-1}知$P_L P_M$是$H$到$\text{Im}(P_L P_M)$的正交投影算子.

下证$\text{Im}(P_L P_M)=L\cap M$.对$\forall P_L P_M x\in \text{Im}(P_L P_M)$,有
\begin{align*}
P_LP_Mx=P_L\left( P_Mx \right) =P_M\left( P_Lx \right) \in L\cap M.
\end{align*}
因此$\text{Im}(P_L P_M)\subset L\cap M$.

再设$y\in L\cap M$,则$P_My=y \in L$,从而$P_L P_My=P_L(P_My)=P_M y=y$.故$y\in \text{Im}P_LP_M$,因此$L\cap M\subset \text{Im}(P_L P_M)$.故$\text{Im}(P_L P_M)= L\cap M$,即$P_L P_M$是$H$到$L\cap M$的正交投影算子.再由\hyperref[theorem:正交投影算子定义]{正交分解的唯一性}知$P_L P_M = P_{L\cap M}$.

\end{enumerate}
\end{proof}

\begin{example}
设 $A \in \mathscr{L}(H)$，称其为正算子，是指
\begin{align*}
(Ax,x) \geqslant 0 \quad (\forall x \in H).
\end{align*}
求证：
\begin{enumerate}[(1)]
\item 正算子必是对称的；
\item 正算子的一切特征值都是非负实数。
\end{enumerate}
\end{example}

\begin{example}
求证：为了 $H$ 的闭线性子空间 $L,M$ 满足 $L \subseteq M$当且仅当 $P_M - P_L$ 是正算子。
\end{example}

\begin{example}
设 $(a_{ij})(i,j=1,2,\cdots)$ 满足 $\sum\limits_{i,j=1}^{\infty} |a_{ij}|^2 < \infty$，在 $l^2$ 空间上，定义映射
\begin{align*}
A: x = \{x_1,x_2,\cdots\} \mapsto y = \{y_1,y_2,\cdots\},
\end{align*}
其中 $y_i \triangleq \sum\limits_{j=1}^{\infty} a_{ij}x_j(i=1,2,\cdots)$。求证：
\begin{enumerate}[(1)]
\item $A$ 是 $H$ 上的紧算子；
\item 又若 $a_{ij} = \overline{a_{ji}}(i,j=1,2,\cdots)$，则 $A$ 是对称紧算子。
\end{enumerate}
\end{example}

\begin{example}
设 $A$ 是 $H$ 上的对称算子，并且存在一组由 $A$ 的特征元组成的 $H$ 的正交规范基。又设
\begin{enumerate}[(1)]
\item $\dim N(\lambda I - A) < \infty \quad (\forall \lambda \in \sigma_p(A)\setminus\{0\})$；
\item $\forall \varepsilon > 0,\ \sigma_p(A)\setminus[-\varepsilon,\varepsilon]$ 只有有限个值。
\end{enumerate}
求证：$A$ 是 $H$ 上的紧算子。
\end{example}


\end{document}